\begin{helper}[Nondefective original-tree confinement]
Let $L$ be a \noderef{ScheideLastExitPath} certificate for $\alpha$.
Suppose that its last inside vertex $v$ belongs to
$X_0=\noderef{TreeSequenceVertexSet}(T)$ and that
$\alpha\notin\noderef{DefectiveColors}(G,\varphi,X_0)$.  Then every
vertex-node of $L$'s path from $y_0$ through $v$ belongs to $X_0$.
Moreover, every edge identity which belongs both to
$\noderef{BoundaryEdges}(G,X)$, where
$X=\noderef{ScheideEnlargedVertexSet}(R)$, and to the full
$(D.\alpha_0,\alpha)$ component through $y_0$ is the crossing edge of
$L$.
\end{helper}
\begin{proof}
The extension field of \noderef{ScheideReducedConfiguration} supplies the
inherited \noderef{ScheideStructuralConfiguration}.  Its globally maximum
field, followed by
\noderef{ScheideGlobalMaximumOrder}, says that $X_0$ is closed under
$\varphi$.  The setup \noderef{ScheideDefectiveSetup} says that
$\alpha_0=D.\alpha_0$ is missing at $y_0\in X_0$.  Thus
\noderef{ColoredVertexSetClosed} forbids every $\alpha_0$-colored
boundary edge of $X_0$.

Let $P,i,v,g$ be the path, last-inside index, last-inside vertex, and
crossing edge recorded by $L$.  The certificate gives $v\in X_0$,
$\varphi(g)=\alpha$, and says that the vertex after $g$ lies outside
$X\supseteq X_0$.  Hence $g$ is an $\alpha$-colored boundary edge of
$X_0$.

Suppose some vertex-node of the segment $y_0Pv$ were outside $X_0$.
Choose the first such vertex-node.  The preceding edge-node $h$ is a
boundary edge of $X_0$.  By \noderef{TwoColorIncidenceGraph}, its color
is $\alpha_0$ or $\alpha$; the preceding closedness argument excludes
$\alpha_0$, so $\varphi(h)=\alpha$.  The path later returns to
$v\in X_0$, and $h\ne g$ because $g$ occurs strictly after $v$ and $P$
is simple.  Thus $h$ and $g$ are two distinct members of
\noderef{BoundaryColorEdges}$(G,\varphi,X_0,\alpha)$.  Expanding
\noderef{DefectiveColors} and \noderef{DefectiveColor}, this says that
$\alpha$ is defective for $X_0$, contrary to the hypothesis.  Therefore
every vertex-node through $v$ lies in $X_0$.

Now let $f$ be any edge identity in both the boundary of $X$ and the full
two-color component.  The support equality in
\noderef{ScheideLastExitPath} puts the edge-node $f$ on $P$.  Its two
neighbors on the full component path are its two endpoint vertex-nodes,
by \noderef{Incident}, \noderef{LooplessMultigraph}, and
\noderef{TwoColorIncidenceGraph}.  All vertex-nodes of $P$ through $v$
lie in $X_0\subseteq X$, while every vertex-node after $v$ lies outside
$X$ by the last-index field of $L$.  Hence the only consecutive pair of
endpoint vertex-nodes separated by the boundary of $X$ occurs across the
edge-node immediately after $v$, namely $g$.  Simplicity of $P$ then
gives $f=g$, proving the second assertion.
\end{proof}
